\section{What it means}

\subsection{The finding, in one sentence}

Deciding which query to run is where this design earns its keep. Holding the
templates and writing the prose are not.

\subsection{Interpretation}

\paragraph{The query-writing gap is real and it survives its fairest
corrections.} Two disjoint question sets put the same gap in the same place.
Quadrupling the token budget did not close it, and gate refusals went up rather
than down under the larger budget. Crediting every failure that belongs to our
own evidence reader still leaves the arm at roughly half the cascade. The
mechanism behind the gap was not isolated. The failure split separates gate
refusals from reader failures, and no experiment attributes the residue to a
specific weakness in the generated queries. \confidence{strong}

\paragraph{The claim that templates hold the accuracy is dead as stated.} It was
described in the claims ledger as the strongest result in the study, and it
rested on a no-templates control that predates the schema-derived capabilities.
Rebuilt, that control matches or beats the cascade on every question set it has
been graded on. The ratio of 3.0 in Table~\ref{tab:seven-arm} prices a missing
capability as much as it prices hand-authoring. \confidence{strong}

\paragraph{What the templates do still hold is less than it first looked.} The
refusal-quality gap in Table~\ref{tab:abstention} was wide on the catalogue and
on the blind set, and it disappeared on both generated sets once the
answer-writing fix landed. The most likely reading is that the earlier gap was a
defect in the generic path rather than a property of templates. That reading was
not tested directly. The one property the templates hold without qualification
is reproducibility: identical answers on repeat runs, so any difference the
cascade shows against another configuration is a real difference.
\confidence{suggestive}

\paragraph{The question set moved the numbers as much as the code did.} The same
system spans 0.72 to 0.98 across the six measurements in
Fig.~\ref{fig:six-numbers}. Two-thirds of what was first attributed to the
question set turned out to be software defects that hit naturally-phrased
questions harder than catalogue questions. A bias measurement taken over a
defective pipeline measures the defect too. \confidence{established}

\paragraph{The top three configurations are not separable.} The prose-writing
arm varies by 0.0283 against itself over identical inputs. Every difference
among the top three arms on either generated set is smaller than that. The order
of the top two also swapped between the two sets. Any table that ranks those
three describes one run of one question set. \confidence{strong}

\subsection{How this fits what was already known}

The size of the query-writing gap is consistent with the published difficulty of
the task. Zero-shot execution accuracy on unseen production graphs averages
0.036 across 11 models \citep{pipecypher}, and a peer-reviewed cross-dialect
benchmark reports 35.40\% \citep{gqlbench}. Our arm scores well above both,
which is what a good-faith build given the schema, the seeds and a repair loop
should do. It is still half the cascade.

Nothing here makes the mechanism novel. Choosing templates over generation is a
named family in knowledge-base question answering \citep{aqqu, tqa, kbqasurvey},
and the modern constrained-decoding line is stronger and better evaluated
\citep{picard}. AgCyRAG independently built this architecture in this domain and
wrote down the same design argument, with no quantitative evaluation of any kind
\citep{agcyrag}. The contribution available to this project is the measurement,
not the design.

The grading protocol is a live vulnerability. Judge-based evaluation in this
area has been shown to swing win rates by more than 30\% from position alone
\citep{gzeng}. Our protocol is better than the one that critique dismantles: a
different model family, given the expected assertions and the full evidence, one
agent per answer. There is still no agreement statistic against a human.

\subsection{Limitations, ranked by what they would change}

\begin{enumerate}
  \item \textbf{No external baseline exists.} Every comparison in this report is
    Ariadne against Ariadne. No result here may be described as beating anything
    outside this repository. This blocks the federation claim entirely, since
    the merged-graph counterfactual has never been built.
  \item \textbf{Both graders are language models, neither validated against a
    human.} The second grader in Table~\ref{tab:cross-grader} emitted two
    distinct scores and cannot resolve differences of a few hundredths. It also
    ran less generous on our grader's weakest arm than on its strongest, which
    is the pattern that would matter if it held. This changes every absolute
    correctness value in the report. It does not close the two-to-one gap, which
    survives the largest offset the second grader produced.
  \item \textbf{Thirty questions per blind set.} A difference of a few
    hundredths is noise at that size, which is exactly the regime the top three
    configurations sit in.
  \item \textbf{One domain, one served model, one span of days.} Nothing here
    generalises to another model without re-running, and the served model's
    quantisation is recorded as unknown rather than guessed.
\end{enumerate}

\subsection{Confidence roll-up}

Templates beat a model writing the query, at roughly two to one, on anchored
questions in this domain. \confidence{strong} The near-perfect catalogue figure
was substantially an artifact of a co-authored question set.
\confidence{established} The templates are not where the correctness lives.
\confidence{strong} Federation beats merging. \confidence{unknown}
